\documentclass[10pt,a4paper]{article}
\usepackage[margin=20mm]{geometry}
\usepackage{amsmath,booktabs,graphicx,hyperref}
\hypersetup{colorlinks=true,urlcolor=blue,citecolor=blue,linkcolor=blue}
\graphicspath{{../figures/airfoil_regimes/}}
\setlength{\parindent}{0pt}
\setlength{\parskip}{5pt}
\title{Airfoil functional reduction:\\geometry neighborhoods and operating regimes}
\author{}
\date{}
\begin{document}
\maketitle
\vspace{-15mm}
\section{Where do similar shapes fall in the functional map?}
The question is whether a two-dimensional aerodynamic representation preserves
geometric neighborhoods, and whether its clusters describe performance across
different operating regimes. The existing population contains 170 NACA
four-digit airfoils: camber $m/c\in\{0,0.02,0.04,0.06,0.08\}$, camber position
$p/c\in\{0.2,0.4,0.6,0.8\}$, and thickness $t/c=0.06,0.08,\ldots,0.24$.
Symmetric sections occur once per thickness, with position recorded as zero.
Cached NeuralFoil predictions at Mach zero cover a $50\times50$ grid on
$\alpha\in[-15,15]^\circ$, $Re\in[10^6,2\times10^6]$.

The saved multivariate functional PCA uses uniform quadrature in $\alpha$ and
linear $Re$, with each coefficient divided by its integrated population
standard deviation $s_c$. This is a discretized product-space functional
metric, following the multivariate functional framework of
\cite{happ2016}. The first two PCs retain 88.29\% of variance; seven retain
at least 99\%. Scores are unwhitened and their saved signs are preserved.

Shape distance uses upper and lower surface coordinates $q_i^u(x),q_i^l(x)$
at 201 cosine-spaced mean-line stations, with normal thickness offsets and
the finite trailing edge of the NACA family \cite{jacobs1933}:
\[
 d_g(i,j)^2=\frac12\int_0^1
 \left(\|q_i^u(x)-q_j^u(x)\|^2+\|q_i^l(x)-q_j^l(x)\|^2\right)\,dx.
\]
This measures registered contour RMS in chord units, rather than standardized
parameter distance or nearest-point distance between curves.

\begin{figure}[h]
\centering\includegraphics[width=\linewidth]{geometry_map.pdf}
\caption{Identical global functional coordinates, colored by geometric parameters.
Gray curves connect increasing thickness at fixed camber and position. Position
zero for symmetric sections is a bookkeeping convention.}
\end{figure}

\begin{center}
\begin{tabular}{lrr}
\toprule Representation & Distance Spearman correlation & Five-neighbor overlap\\
\midrule Two global PCs & 0.673 & 36.2\%\\
Seven global PCs & 0.666 & 44.7\%\\
Complete fields & 0.666 & 45.3\%\\
Random neighborhood expectation & --- & 3.0\%\\
\bottomrule
\end{tabular}
\end{center}
Overlap averages the fraction of each shape neighborhood recovered by the
aerodynamic neighborhood. Correlations use all 14,365 unordered pairs;
these are dependent observations, so no significance test is asserted.
PC1 correlates with camber ($\rho=0.797$), while PC2 mixes camber
($-0.581$), thickness ($-0.550$), and position ($-0.398$).
Geometric proximity is partly preserved; the improvement from two to seven
PCs concerns local neighborhoods even though global distance correlation barely changes.

\newpage
\section{Partitioning the operating space and defining performance}
Define angle-magnitude bins low $[0,5]$, middle $[5,10]$, high $[10,15]$
degrees, and Reynolds bins low $[1,4/3]$, middle $[4/3,5/3]$,
high $[5/3,2]$, in millions. Their Cartesian product gives nine regimes
$B_r$, named angle--Reynolds. Boundaries have zero measure. Angle bins combine
negative and positive incidence with equal weight; they are chosen operating
intervals, not identified attached-flow or stall regimes.

For any quantity $f$, let
\[
 \langle f\rangle_r=\frac{1}{|B_r|}\int_{B_r} f(\alpha,Re)\,d\alpha\,dRe,
 \qquad
 P_{ir}=\left(\langle C_D\rangle_r,\langle C_L\rangle_r,
 \langle|C_L|\rangle_r,\left\langle\frac{|C_L|}{C_D}\right\rangle_r\right).
\]
Lower mean drag, higher lift magnitude, and higher efficiency are distinct
performance objectives in every regime. At low angle, drag and useful lift
must be examined jointly; at middle angle, efficiency is the primary screening
quantity; at high angle, drag penalties and maintained lift magnitude matter.
There is no single optimum without a lift requirement and mission weighting.
Absolute lift permits bidirectional screening; it is not a positive-lift
mission objective. Signed mean lift is saved separately to expose camber bias.
Efficiency averages pointwise ratios, rather than dividing mean lift by mean drag.

The computation integrates piecewise-linear nodal quantities across clipped
bin boundaries. This avoids assigning an entire grid cell to one regime;
volume-weighted bin quadratures recover the original whole-domain quadrature.
The interpolation is applied after evaluating absolute lift and efficiency at
the nodes. No aerodynamic solver or new confidence estimate is introduced.

We cluster the 27 descriptors
$(\langle C_D\rangle_r,\langle|C_L|\rangle_r,
\langle|C_L|/C_D\rangle_r)_{r=1}^9$ after standardizing each descriptor across
the population. This gives equal marginal weight to each regime and metric.
These correlated metrics deliberately represent several performance objectives;
they are not independent physical coordinates. K-means uses $k=4$, seed 42,
20 initializations and at most 500 iterations. Cluster numbers R1--R4 are
ordered by mean camber, then thickness. Four clusters retain continuity with
the previous experiment; an optimal cluster count is not established.

\begin{figure}[h]
\centering\includegraphics[width=\linewidth]{regime_clusters.pdf}
\caption{The same regime-performance assignments shown on the existing global
functional map and on a PCA of the 27 standardized descriptors. Clustering
uses all descriptors; the right-hand display retains 73.73\% of their variance.}
\end{figure}

The descriptor-space silhouette is 0.323 and relative centroid SSE is 0.377.
Agreement with global seven-PC clusters is ARI $=0.488$, compared with
ARI $=0.980$ between the previous two- and seven-PC partitions.
Thus the definition of similarity changes substantially when performance
summaries receive equal weight across operating regimes.

\newpage
\section{Performance tradeoffs between regime clusters}
\begin{figure}[h]
\centering\includegraphics[width=\linewidth]{performance.pdf}
\caption{Cluster means of per-airfoil regime averages. Columns order angle
magnitude first, then Reynolds number. Color scales differ between metrics.
Member medians and 10--90\% quantiles are saved in the companion tables;
the means do not imply that every member shares the same ranking.}
\end{figure}

\begin{center}\small
\begin{tabular}{lrrrl}
\toprule
Cluster & $n$ & Mean $m/c$ (\%) & Mean $t/c$ (\%) & Interpretation\\
\midrule
R1 & 55 & 2.22 & 16.95 & Low drag; strongest high-angle efficiency\\
R2 & 54 & 5.15 & 16.56 & Intermediate camber and efficiency\\
R3 & 27 & 5.48 & 7.26 & Thin; high-angle drag penalty\\
R4 & 34 & 7.41 & 15.53 & High camber; strongest low-angle efficiency\\
\bottomrule
\end{tabular}

\end{center}

R4 has the largest mean low-angle efficiency: 96.2, 102.2, and 106.6 across
increasing Reynolds bins. R1 instead has the largest mean middle-angle
efficiency (70.8--78.8) and high-angle efficiency (60.2--72.0).
At high angle and low Reynolds, R1 has mean drag 0.022 and efficiency 60.2,
whereas the thinner R3 group has mean drag 0.107 and efficiency 20.2.
These observations describe this surrogate and sampled family; thickness
alone is not established as the cause.

Within-cluster variability remains material. For high-angle/low-Reynolds
efficiency, member 10--90\% ranges are 51.7--66.6 for R1,
28.4--53.9 for R2, 7.4--41.0 for R3, and 22.0--40.7 for R4.
These are descriptive population spreads, not confidence intervals.
In particular, a cluster label does not determine the best airfoil for a
specific operating point or prescribed lift.

\newpage
\section{Regime-local functional structure}
For each regime, a second K-means fit uses the complete fields restricted to
that regime, with the original global coefficient scales held fixed. All
partitions are compared using the same regime-specific field metric:
\[
 \mathrm{SSE}_r=\frac{\sum_i\|F_i-\bar F_{g(i)}\|_r^2}
 {\sum_i\|F_i-\bar F\|_r^2},\qquad
 \|F\|_r^2=\sum_{c\in\{D,L\}}\frac{\langle F_c^2\rangle_r}{s_c^2}.
\]
The reconstruction diagnostic $E_j$ replaces the cluster mean in the
numerator by the saved $j$-PC reconstruction. Local $V_2$ is the fraction of
regime field variance captured by a fresh two-dimensional PCA under this
same restricted metric. These are in-population diagnostics, not held-out errors.

\begin{center}\small
\begin{tabular}{lrrrrrr}
\toprule
Regime & $E_2$ & $E_7$ & Local $V_2$ & Global SSE & Local SSE & ARI\\
\midrule
low-low & 0.034 & 0.008 & 0.990 & 0.290 & 0.101 & 0.360\\
low-middle & 0.033 & 0.008 & 0.992 & 0.288 & 0.095 & 0.360\\
low-high & 0.032 & 0.008 & 0.994 & 0.287 & 0.092 & 0.362\\
middle-low & 0.117 & 0.014 & 0.936 & 0.302 & 0.203 & 0.460\\
middle-middle & 0.119 & 0.014 & 0.947 & 0.319 & 0.186 & 0.397\\
middle-high & 0.117 & 0.016 & 0.953 & 0.326 & 0.179 & 0.383\\
high-low & 0.177 & 0.008 & 0.856 & 0.269 & 0.179 & 0.664\\
high-middle & 0.135 & 0.006 & 0.876 & 0.217 & 0.172 & 0.751\\
high-high & 0.118 & 0.008 & 0.893 & 0.202 & 0.174 & 0.807\\
\bottomrule
\end{tabular}

\end{center}
Global denotes the seven-PC whole-domain partition. ARI compares that
partition with the regime-local partition. Local fits optimize their own
in-population SSE, so their lower SSE is a descriptive baseline rather than
evidence of predictive superiority.

\begin{figure}[h]
\centering\includegraphics[width=.67\linewidth]{local_partitions.pdf}
\caption{Regime-local cluster memberships drawn on the same global 2D points.
Rows increase angle magnitude and columns increase Reynolds number. Colors
are local to each panel and cluster identifiers are not matched across panels.}
\end{figure}

Two global PCs lose only 3.2--3.4\% of low-angle field energy, but
11.8--17.7\% at high angle. Two locally fitted PCs retain 85.6--89.3\%
at high angle and 99.0--99.4\% at low angle. Seven global PCs keep regime
errors below 1.6\%. Low-angle local partitions have ARI near 0.36 with the
global partition; high-angle agreement rises from 0.66 to 0.81 across Reynolds
bins. A global map can therefore obscure regime-specific structure even
when the global four-cluster partition barely changes between two and seven PCs.

\newpage
\section{Counterexamples, sensitivity, and scope}
\begin{figure}[h]
\centering\includegraphics[width=.80\linewidth]{contour_examples.pdf}
\caption{Deterministically selected counterexamples. Top: the most separated
2D pair among the closest 5\% of contour pairs. Bottom: the most different
contour pair among the closest 5\% of 2D pairs. The registered RMS definition
and all pair distances are saved; these examples are illustrative extremes.}
\end{figure}
NACA4208/NACA4610 have contour RMS 0.0136 chord (3.8th distance percentile)
but 2D distance at the 64.8th percentile. NACA4810/NACA6224 have contour RMS
0.0576 chord (86.4th percentile) while their 2D distance is at the 3.2nd
percentile. Close projected responses can correspond to substantially
different shapes, and similar shapes need not have close projected responses.

The regime-performance partition is sensitive to scientific choices.
Positive-only and negative-only incidence give ARI 0.381 and 0.500 against
the bidirectional fit. Moving angle boundaries to $(0,4,9,15)^\circ$ or
$(0,6,11,15)^\circ$ gives ARI 0.460 and 0.515. Omitting one middle-angle
regime gives ARI 0.446--0.450; omitting one low-angle regime gives 0.981.
Ten 80\% subsample fits, each refitting the standardization and predicting
the whole population, yield mean ARI 0.899, minimum 0.460. Using three or
five clusters gives ARI 0.478 or 0.950 against the four-cluster fit.
These assignments are useful exploratory performance groups, not stable
physical classes established independently of the chosen bins and metrics.

The data are model predictions, without retained pointwise confidence,
CFD or experimental validation. NeuralFoil exposes confidence diagnostics
\cite{neuralfoil}, but this cache cannot support confidence-conditioned claims.
There is no identified stall boundary, trim constraint, fixed-lift drag
comparison, extrapolation test, or optimized mission weighting. The next useful
check is a prescribed-lift, positive-incidence study with confidence filtering
and higher-fidelity validation of representative members.

\textbf{Reproduction.} Run \texttt{python notebooks/airfoil\_regimes.py} from
the repository root with the numerical analysis dependencies installed.
The script verifies the cache hash and saved scores, regenerates CSV/JSON
evidence and vector figures, and updates the report tables. From
\texttt{report/}, run \texttt{pdflatex airfoil\_regimes},
\texttt{bibtex airfoil\_regimes}, then \texttt{pdflatex airfoil\_regimes} twice.
Software versions, source hashes, seed and generation settings are recorded
in \texttt{data/airfoil\_regimes/run\_metadata.json}.

\small
\bibliographystyle{plain}
\bibliography{references}
\end{document}
